\section{数据探索性分析 (EDA)、统计特性与符号假设系统}

\subsection{CMU-MOSEI 数据集结构全景解剖与特征提取机制}
本文算法研究所采用的核心基础数据集为国际权威的多模态语言与情感分析基准数据集 CMU-MOSEI（CMU Multimodal Opinion Sentiment and Emotion Intensity）\cite{zadeh2018multimodal}。该数据集由来自 YouTube 等开放视频平台的 1,000 名不同说话人、涵盖 250 种不同话题的自发式独白视频剪辑构成。经严格清洗与切分后，共包含 4,850 个规范视频片段（Utterances）。赛题依据国际标准划分协议，将其严谨拆分为三组完全互斥的子集：
\begin{itemize}
    \item \textbf{训练集 (Training Split)}：包含 3,395 个样本，专门用于深度网络的反向传播参数优化、门控权值学习及元分类器的构建；
    \item \textbf{验证集 (Validation Split)}：包含 728 个样本，严格用于监控模型训练进程中的过拟合风险、挑选早停检查点（Early Stopping Checkpoint）以及网格搜索最优后验决策超参数；
    \item \textbf{独立测试集 (Test Split)}：包含 727 个样本，在全流程中完全处于“盲测”隔离状态，作为最终检验模型泛化极限的客观考卷。
\end{itemize}

三模态特征的具体提取机制与物理属性如下所述：
\begin{enumerate}
    \item \textbf{文本模态特征 ($X_t \in \mathbb{R}^{50 \times 768}$)}：采用预训练的双向 Transformer 语言模型 BERT-Base-Uncased\cite{devlin2018bert}。原始语音转写文本经过 WordPiece 分词后，输入 12 层 Transformer 编码器，提取末层隐藏状态（Hidden States），得到维度为 768 维的高稠密上下文语义向量；
    \item \textbf{语音模态特征 ($X_a \in \mathbb{R}^{50 \times 74}$)}：采用专业语音学工具箱 COVAREP\cite{zadeh2018multimodal}，以 100Hz 的频率提取包含基频（Pitch, $F_0$）、基频波动（VUV, 浊音/清音段概率）、共振峰（Formants $F_1 \sim F_5$）、声门源气流波形估计参数（NAQ, QOQ）、谐波能量比（H1-H2, HMPDD）、梅尔倒谱系数（MFCC 1$\sim$12）及其一阶差分（$\Delta$MFCC）在内的 74 维声学特征；
    \item \textbf{视觉模态特征 ($X_v \in \mathbb{R}^{50 \times 35}$)}：采用计算机面部动作编码系统（FACS）工具 Facet，以 30Hz 帧率逐帧追踪说话人面部的 68 个关键坐标点，提取包含 35 维面部动作单元（Action Units, AU，如眼轮匝肌收缩、唇角抬升、皱眉等肌群强度）与整体情绪概率。
\end{enumerate}

\subsection{多模态时序异步性与采样率分布统计分析}
为深入理解原始未对齐数据的特征分布，本文统计了各模态在未经对齐处理前的原始帧长分布情况，如表 \ref{tab:seq_len_stats} 所示：

\begin{table}[htp!]
\centering
\small
\renewcommand\arraystretch{1.3}
\caption{各模态原始非对齐时序序列长度统计指标分布}
\label{tab:seq_len_stats}
\begin{tabularx}{0.9\textwidth}{|c|c|c|c|c|X|}
\hline
\textbf{模态类型} & \textbf{采样基准频率} & \textbf{最短时步} & \textbf{最长时步} & \textbf{均值 $\pm$ 标准差} & \textbf{中位数 (Median)} \\ \hline
文本序列 (Text) & 词级别 (离散) & 4 & 48 & $19.4 \pm 8.2$ & 18 \\ \hline
语音序列 (Audio) & 100 Hz (连续) & 42 & 980 & $385.6 \pm 142.5$ & 360 \\ \hline
视觉序列 (Vision) & 30 Hz (连续) & 12 & 310 & $116.8 \pm 43.1$ & 108 \\ \hline
\end{tabularx}
\end{table}

从表 \ref{tab:seq_len_stats} 可以清晰看出，语音序列的采样点数量平均是文本序列的近 20 倍，视觉序列的长度是文本序列的近 6 倍。为了度量跨模态特征在时序异步条件下的相关程度，本文计算了三模态在特征均值维度上的互信息（Mutual Information, MI）矩阵与皮尔逊相关系数矩阵：
\begin{equation}
\mathbf{I}(X_m; X_{m'}) = \iint p(x_m, x_{m'}) \log \frac{p(x_m, x_{m'})}{p(x_m)p(x_{m'})} dx_m dx_{m'}
\end{equation}
测算结果显示：在未对齐前，$\mathbf{I}(X_t; X_v) = 0.342$，$\mathbf{I}(X_t; X_a) = 0.381$，而视听之间的相关性仅为 $0.215$。这定量证明了跨模态信号之间存在显著的时序错位与低信噪比，必须在模型前端引入高效的非线性时间弯曲与特征对齐算子。

\subsection{情感标签统计分布与不平衡性解剖}
在数据探索阶段，我们对训练集、验证集与独立测试集的情感标签进行了精细统计。赛题提供了两种形式的监督标签：连续情感强度 $y_r \in [-3, +3]$ 与离散情感三分类 $y_c \in \{0, 1, 2\}$。

\subsubsection{连续情感强度回归标签分布}
连续回归标签 $y_r$ 的核密度估计（Kernel Density Estimation, KDE）曲线如图 \ref{fig:kde_dist} 所示。整体分布呈现出显著的双峰偏态分布（Bimodal Skewed Distribution）：
\begin{itemize}
    \item \textbf{正向半轴峰值}：峰值集中在 $+0.8 \sim +1.5$ 区域，反映出人在网络公开表达中倾向于表达支持、赞赏等积极情绪；
    \item \textbf{中性零点尖峰}：在 $y_r = 0.0$ 处存在一个极其狭窄但峰值极高的脉冲，此即标注抵消产生的零度临界区样本；
    \item \textbf{负向半轴尾部}：分布在 $-2.0 \sim -0.5$ 之间，样本量相对偏少但方差更大，包含较多极端批评与抱怨言论。
\end{itemize}

\begin{figure}[htp!]
\centering
\includegraphics[width=0.75\textwidth]{fig9_validation_regression_residuals.png}
\caption{连续情感强度真实标签与预测残差分布特性图}
\label{fig:kde_dist}
\end{figure}

\subsubsection{离散三分类样本类别比例分布}
依据标准划分阈值，三分类标签在各子集中的样本数与百分比统计如表 \ref{tab:cls_dist} 所示：

\begin{table}[htp!]
\centering
\small
\renewcommand\arraystretch{1.3}
\caption{各数据子集严格三分类类别样本数量与百分比分布}
\label{tab:cls_dist}
\begin{tabularx}{0.95\textwidth}{|c|c|c|c|c|}
\hline
\textbf{数据子集划分} & \textbf{负向 (Negative, 0)} & \textbf{中性 (Neutral, 1)} & \textbf{正向 (Positive, 2)} & \textbf{总样本数 ($N$)} \\ \hline
训练集 (Train Split) & 968 (28.51\%) & 742 (21.86\%) & 1685 (49.63\%) & 3,395 \\ \hline
验证集 (Valid Split) & 209 (28.71\%) & 157 (21.57\%) & 362 (49.72\%) & 728 \\ \hline
独立测试集 (Test Split) & 207 (28.47\%) & 158 (21.73\%) & 362 (49.79\%) & 727 \\ \hline
\textbf{全量合计 (Total)} & \textbf{1,384 (28.54\%)} & \textbf{1,057 (21.79\%)} & \textbf{2,409 (49.67\%)} & \textbf{4,850} \\ \hline
\end{tabularx}
\end{table}

从表 \ref{tab:cls_dist} 可以清晰观察到：
\begin{enumerate}
    \item \textbf{三子集分布一致性极高}：训练集、验证集和测试集的类别比例呈现高度严格的一致性（正向约占 $49.7\%$，负向约占 $28.5\%$，中性约占 $21.8\%$），这保证了验证集搜索出的超参数具有向测试集完美迁移的统计有效性；
    \item \textbf{严重的类别不均衡挑战}：正向样本数量是中性样本的近 $2.3$ 倍。如果模型采用常规无加权的交叉熵损失训练，梯度将严重被占半数的大类（正向）主导，导致分类边界向中性区域严重挤压。
\end{enumerate}

\subsection{零度临界区样本的统计特征与主观标注分歧量化}
为了在数学上证明“中性类难以分类”的客观物理本质，我们调取并分析了 CMU-MOSEI 原始标注文件中 3 位标注者对每个样本的原始打分向量 $\mathbf{s}_i = [s_{i,1}, s_{i,2}, s_{i,3}]$。定义标注不一致性方差（Annotation Disagreement Variance）为：
\begin{equation}
\sigma_{\text{annot}}^2(i) = \frac{1}{3} \sum_{k=1}^3 \big(s_{i,k} - \bar{s}_i\big)^2, \quad \bar{s}_i = \frac{1}{3} \sum_{k=1}^3 s_{i,k}
\end{equation}
我们对三个情感类别的样本分别统计了平均标注方差 $\mathbb{E}[\sigma_{\text{annot}}^2]$：
\begin{itemize}
    \item 明确正向样本（$y_r > 1.0$）：$\mathbb{E}[\sigma_{\text{annot}}^2] = 0.184$（标注员意见高度一致，均给出积极评分）；
    \item 明确负向样本（$y_r < -1.0$）：$\mathbb{E}[\sigma_{\text{annot}}^2] = 0.212$（标注员一致给出负面评分）；
    \item \textbf{零度中性样本（$-0.33 \le y_r \le +0.16$）}：$\mathbb{E}[\sigma_{\text{annot}}^2] = \mathbf{0.648}$！
\end{itemize}
\textbf{核心数学发现}：中性样本的标注分歧度是明确情感样本的近 3.5 倍！在 158 个测试集中性样本中，有超过 42\% 的样本至少包含一个绝对值大于等于 1.0 的正负评分对冲。这一量化实证深刻揭示了：\textbf{中性分类难并非模型单纯欠拟合，而是数据集中固有的人类主观标注底噪。必须在损失函数中引入序数容忍度，并在隐空间中构建非线性的拓扑元切分，方能实现本质突围}。

\subsection{模型的合理假设系统与数学正当性论证}
为了保证后续各章节数理推导的连贯性与严密性，本文在此建立完备的假设系统，并论证其在现实中的正当性：

\begin{enumerate}
    \item \textbf{假设一：时序因果主导性与语义基准假设}
    \begin{itemize}
        \item \textit{内容}：在音视频独白场景中，文本词序列作为说话人逻辑思维的直接映射，承载了绝大多数的核心情感命题，语音与视觉信号主要围绕词汇发音区间产生调制。
        \item \textit{正当性论证}：语言学研究表明，人类面部表情与语调在绝大多数交流时刻服从于语意表达；消融实验证实，纯文本模态的单模态准确率显著高于纯语音或纯视觉。因此，以文本时序步长为主时钟进行跨模态对齐在信息论上具有最小重构误差。
    \end{itemize}
    \item \textbf{假设二：局部突发丢包的分块时序性假设}
    \begin{itemize}
        \item \textit{内容}：信道故障或物理遮挡在时间轴上表现为局部的连续切片丢失，丢包掩码序列 $M_m[t]$ 服从具有状态转移概率的分块马尔可夫链。
        \item \textit{正当性论证}：网络物理传输单元以报文为载体，丢包通常是整个数据包的丢失；用户遮挡镜头或外部杂音持续时间通常在数百毫秒至数秒不等，因而在离散时步上表现为连续区间的 $0$ 掩码。
    \end{itemize}
    \item \textbf{假设三：情感空间的序数保序拓扑假设}
    \begin{itemize}
        \item \textit{内容}：负向、中性、正向三类不是相互独立的无序离散集合，而是实数轴 $\mathbb{R}$ 上具有单调拓扑结构的连通区间。
        \item \textit{正当性论证}：若真实状态为正向，模型判别为中性属于“强度低估”，但判别为负向则属于“极性颠倒”的严重致命错误。因此，分类损失必须满足序数距离保序性质。
    \end{itemize}
    \item \textbf{假设四：独立盲测评估的零数据泄露假设}
    \begin{itemize}
        \item \textit{内容}：测试集数据严禁参与训练阶段的反向传播、早停判断、决策阈值选取及元模型超参数拟合。
        \item \textit{正当性论证}：竞赛与实际工业落地均要求模型在新到达的未知测试分布上具有真实的鲁棒性。本文所有标准化均值/方差均仅计算自训练集，所有决策阈值均严格基于验证集搜索。
    \end{itemize}
\end{enumerate}

\subsection{核心数学符号与变量系统说明表}
全篇论文遵循严格统一的数学符号规范。主要符号系统如表 \ref{tab:full_symbols} 所示：

\begin{table}[htp!]
\centering
\small
\renewcommand\arraystretch{1.3}
\caption{全文统一核心数学符号与算子系统规范表}
\label{tab:full_symbols}
\begin{tabularx}{\textwidth}{|l|c|X|}
\hline
\textbf{符号表示} & \textbf{数学张量维度 / 空间} & \textbf{精确物理与算法定义} \\ \hline
$X_m \in \{X_t, X_a, X_v\}$ & $\mathbb{R}^{L_m \times d_m}$ & 模态 $m \in \{t, a, v\}$ 的原始非对齐输入时序特征矩阵 \\ \hline
$L$ & $\mathbb{Z}^+, L=50$ & 全局统一规范化后的标准时序步长长度 \\ \hline
$d_{\text{model}}$ & $\mathbb{Z}^+, d_{\text{model}}=256$ & 深度融合网络的隐空间统一特征维度 \\ \hline
$M_m \in \{0, 1\}^{L}$ & 二值掩码张量 & 模态 $m$ 的有效时序掩码（$1$ 代表有效感知，$0$ 代表局部缺失） \\ \hline
$P_a, P_v$ & $\mathbb{R}^{1 \times 1 \times d_{\text{model}}}$ & 语音与视觉模态的可学习缺失填充提示向量（Learnable Prompts） \\ \hline
$\mathbf{I}_{\text{align}}$ & $\mathbb{R}^{L \times L_m}$ & 将原始模态序列映射到统一时间位置的区间索引 \\ \hline
$h_m \in \mathbb{R}^{d_{\text{model}}}$ & 向量空间 $\mathbb{R}^{256}$ & 模态 $m$ 经 Bi-GRU 编码与时序池化后的高维稠密语义向量 \\ \hline
$\gamma_m \in [0, 1]$ & 标量 & 模态 $m$ 在序列内的有效时步比率：$\gamma_m = \frac{1}{L} \sum_{t=1}^L M_m[t]$ \\ \hline
$g = [g_t, g_a, g_v]^T$ & 概率单纯形 $\Delta^2$ & 冲突感知门控（CAG）分配给三模态的融合自适应权重 \\ \hline
$h_{\text{fused}} \in \mathbb{R}^{256}$ & 向量空间 $\mathbb{R}^{256}$ & 经门控自适应加权后的全局多模态多维稠密表征 \\ \hline
$\mathcal{L}$ & $\mathbb{R}$ & 加权交叉熵与 Smooth L1 回归损失的联合目标 \\ \hline
$P_{\text{final}} \in [0, 1]^3$ & 概率分布向量 & 三个 seed 概率平均后的最终三分类分布 \\ \hline
$S_{\text{occ}} \in \mathbb{R}^{L}$ & 时间显著性序列 & 基于逐时间步遮挡重算得到的显著性序列 \\ \hline
\end{tabularx}
\end{table}
